\section{The Pain of Loss: Exile from the Face of God}

Most imaginations, when they hear the word \textit{hell}, jump immediately to flame. The Fathers and Doctors of the Church begin somewhere deeper, and far more terrible. St.\ John Chrysostom, preaching in Antioch on the Gospel of Matthew, sets the pain of separation above every sensory terror:

\begin{quote}
No doubt hell, and that punishment, is a thing not to be borne. Yet though one suppose ten thousand hells, he will utter nothing like what it will be to fail of that blessed glory, to be hated of Christ, to hear ``I know you not,'' to be accused for not feeding Him when we saw Him an hungered. Yea, better surely to endure a thousand thunderbolts, than to see that face of mildness turning away from us, and that eye of peace not enduring to look upon us.\endnote{St.\ John Chrysostom, \textit{Homilies on the Gospel of St.\ Matthew} 23.9, trans.\ George Prevost and M.\ B.\ Riddle, NPNF1 vol.\ 10.}
\end{quote}

Chrysostom is stating what Catholic theology names the \textit{poena damni}, the pain of loss, and judges unanimously to be the heart and substance of damnation.

The earliest systematic articulation of this truth belongs to St.\ Irenaeus of Lyons in the late second century. Writing in Book V of \textit{Against Heresies}, Irenaeus refuses to picture hell as an external mechanism bolted onto a neutral universe. Everything begins with what communion with God actually is:

\begin{quote}
And to as many as continue in their love towards God, does He grant communion with Him. But communion with God is life and light, and the enjoyment of all the benefits which He has in store. But on as many as, according to their own choice, depart from God, He inflicts that separation from Himself which they have chosen of their own accord. But separation from God is death, and separation from light is darkness; and separation from God consists in the loss of all the benefits which He has in store. \dots\ Now, good things are eternal and without end with God, and therefore the loss of these is also eternal and never-ending. It is in this matter just as occurs in the case of a flood of light: those who have blinded themselves, or have been blinded by others, are for ever deprived of the enjoyment of light. It is not, however, that the light has inflicted upon them the penalty of blindness, but it is that the blindness itself has brought calamity upon them.\endnote{St.\ Irenaeus of Lyons, \textit{Against Heresies} V.27.2, trans.\ Alexander Roberts and W.\ H.\ Rambaut, ANF vol.\ 1.}
\end{quote}

Because God is the source of being, light, truth, and joy, to sever oneself from God by settled apostasy is not merely to forfeit a prize while retaining a quiet neutral existence elsewhere; it is to lose every good that has its source in Him, and because those goods are eternal, the destitution is eternal as well.

Two centuries later, St.\ Gregory of Nazianzus, standing in the pulpit at Nazianzus after plague and hail had ruined the harvest of 373, lifts his people's eyes from temporal calamity to the last tribunal---a judgment-seat ``more just even than fearful; or rather more fearful because it is also just.''\endnote{St.\ Gregory of Nazianzus, \textit{Oration} 16.9, trans.\ Charles Gordon Browne and James Edward Swallow, NPNF2 vol.\ 7.} Heaven, Gregory says, consists solely and above all else in ``the unspeakable light and the vision of the holy and royal Trinity, Which now shines upon them with greater brilliancy and purity and unites Itself wholly to the whole soul.'' And what is its opposite? ``The others among other torments, but above and before them all must endure the being outcast from God, and the shame of conscience which has no limit.''

In the Latin West, St.\ Augustine sets the same foundation in his \textit{Enchiridion}. Following his reply to those who imagine that God's mercy might bring an intermission to the torments of the lost (\S 112), Augustine defines that perpetual death (\textit{mors illa perpetua}) as ``alienation from the life of God'' (\textit{alienatio a vita Dei}), which abides without end and is common to all the damned, whatever gradations may attend their other punishments.\endnote{St.\ Augustine, \textit{Enchiridion} 113 (ch.\ 29, \S 113; following the discussion of the \textit{misericordes} in \S 112).}

St.\ Thomas Aquinas supplies the metaphysical reason why this pain of loss surpasses every other suffering and why it is, in a strict sense, infinite. In the \textit{Prima Secundae} of the \textit{Summa Theologiae} (I-II, q.\ 87, a.\ 4), Aquinas asks whether sin incurs a debt of punishment infinite in quantity. He answers by distinguishing the two movements inside every mortal sin, the very two evils that the prophet Jeremias lamented---``They have forsaken me, the fountain of living water, and have digged to themselves cisterns, broken cisterns, that can hold no water'' (Jer 2:13):

\begin{quote}
Punishment is proportionate to sin. Now sin comprises two things. First, there is the turning away from the immutable good, which is infinite, wherefore, in this respect, sin is infinite. Secondly, there is the inordinate turning to mutable good. In this respect sin is finite, both because the mutable good itself is finite, and because the movement of turning towards it is finite, since the acts of a creature cannot be infinite. Accordingly, in so far as sin consists in turning away from something, its corresponding punishment is the \textit{pain of loss}, which also is infinite, because it is the loss of the infinite good, i.e.\ God. But in so far as sin turns inordinately to something, its corresponding punishment is the \textit{pain of sense}, which is also finite.\endnote{St.\ Thomas Aquinas, \textit{Summa Theologiae} I-II, q.\ 87, a.\ 4, trans.\ Fathers of the English Dominican Province.}
\end{quote}

Mortal sin is a free creature's deliberate aversion (\textit{aversio}) from the uncreated Good for the sake of a created trifle (\textit{conversio ad creaturam}). The pain of sense answers to the finite creature chosen in God's place, and therefore remains finite and varies by degree from sinner to sinner. The pain of loss answers to the infinite God rejected, and therefore takes on an infinity of its own: not that a finite soul can suffer an infinite intensity of sensation, which created nature could not sustain, but that the object forfeited is infinite.

The \textit{Roman Catechism} writes this theology directly into its exposition of Christ's sentence in Matthew 25:41. Every clause of the verdict has its own weight, and the very first three words carry the worst blow:

\begin{quote}
These first words, ``depart from me,'' express the heaviest punishment with which the wicked shall be visited---their eternal banishment from the sight of God, unrelieved by one consolatory hope of recovering so great a good. This divines call ``the pain of loss,'' because in hell, the wicked shall be deprived of the light of the vision of God. The words ``ye cursed,'' which are added, must augment to an extreme degree, their wretched and calamitous condition. If when banished from the Divine presence, they could hope for blessing of any sort, it might be to them some source of consolation; but deprived of every such expectation that could alleviate calamity, the divine justice, whose severity their crimes have provoked, pursues them with every species of malediction.\endnote{\textit{Catechism of the Council of Trent}, Part I, Creed Article VII, trans.\ Donovan, p.\ 64.}
\end{quote}

Why, then, does the pain of loss frighten sinners so little in the present life? Because while we live in the body, surrounded by food, company, noise, work, and distraction, a man can spend years in mortal sin without feeling the emptiness of living without God. Created goods still mask the void; the broken cisterns still hold enough muddy water to dull the thirst. Both St.\ Francis de Sales and St.\ Alphonsus Liguori pierce that illusion using a scene from the Second Book of Kings. When Absalom was allowed to return from exile in Gessur to Jerusalem, King David refused to let his son enter the royal presence; after two years Absalom cried out that it would have been better to stay in exile, saying, ``I beseech thee therefore that I may see the face of the king: and if he be mindful of my iniquity, let him kill me'' (2 Kings 14:32). St.\ Francis de Sales turns the scene directly upon the soul: ``Now, if Absalom found the privation of the amiable face of his father, David, more grievous to him than his banishment, good God! what grief will it cause to be forever excluded from the sight of thy most sweet and gracious countenance?''\endnote{St.\ Francis de Sales, \textit{Introduction to the Devout Life} I.15 (Seventh Meditation: On Hell), Pustet ed., pp.\ 38--39.}

St.\ Alphonsus Liguori deepens the point with a piercing observation on the moment of death. Here on earth, he writes, sinners live in darkness and ``are content to live for months and years without God,'' scarcely grasping what His absence means. Yet the instant the soul separates from the body and all earthly distractions drop away, she perceives with unbearable clarity that she was made for God alone:

\begin{quote}
Hell does not consist in the darkness, stench, shrieks and fire; the pain which constitutes Hell, is the pain of having lost God. \dots\ At her separation from the body, the soul is naturally drawn to God, but sin drags her away, and sends her to Hell, at a distance from him. ``Your iniquities,'' says the Prophet Isaias, ``have divided between you and your God.'' The entire, then, of the Hell of the damned consists in that first word of the sentence of their condemnation,---``Depart from me, you cursed.''\endnote{St.\ Alphonsus Maria de Liguori, \textit{Preparation for Death}, Consideration XXVI, Third Point, trans.\ by a Catholic clergyman (Dublin: James Duffy, 1843), pp.\ 323--324.}
\end{quote}

There lies the deepest torment of the lost soul: she does not cease to be a rational creature made for the infinite Good, yet by her own fixed hatred she refuses the very Good toward which her whole nature was fashioned. As St.\ Paul describes it in his second letter to the Thessalonians, it is ``eternal punishment in destruction, from the face of the Lord and from the glory of his power'' (2 Thess 1:9)---banishment from the only Face that could ever satisfy the human heart.
